% GENERATED FILE. Edit canonical lessons, then run npm run generate:books.
% canonical-source-hash: fnv1a64:8ac73b614eea0475
% canonical-lessons: LA-C217-tegula, LA-C217-fornax, LA-C217-lectica, LA-C217-vehiculum, LA-C217-currus

\chapter{Roof tile, Furnace, Litter, Vehicle, Chariot}
\label{ch:la-roof-tile-furnace-litter-vehicle-chariot}
\hlchaptermodality{217}

\begin{chapteropening}
\textbf{By the end of this chapter:} \emph{I can say a roof tile, a furnace, a litter, a vehicle and a chariot.}

Complete the last lesson of chapter 217: I can say a roof tile, a furnace, a litter, a vehicle and a chariot.
\end{chapteropening}

\section[tēgula, tēgulae]{tēgula, tēgulae — a roof tile}

\label{lesson:LA-C217-tegula}



\begin{quote}
Before the new one: say the Latin for worn out, then the Latin for wakeful.
\end{quote}

\subsection*{You'll want to know: tēgula, tēgulae}
\textbf{tēgula, tēgulae} — \textquotedblleft{}a roof tile\textquotedblright{}.

\begin{cousinweb}
\textbf{tēgula} is built on \emph{tegō}, \textquotedblleft{}to cover\textquotedblright{} (EtymDB 2.1), with \emph{-ula}, an ending for a tool: the thing that covers. English \emph{tile} was borrowed from it early.
\end{cousinweb}

\begin{grammarlens}[title={What you've built}]
One more word for this chapter.
\end{grammarlens}

\subsection*{Your turn}
\begin{itemize}
\raggedright
  \item \emph{Say it:} \emph{tēgula, tēgulae}
  \item \emph{Say it:} \emph{tēgula, tēgulae}, once more
  \item \emph{Say it:} say \emph{tēgula, tēgulae}
  \item \emph{Recall:} say the Latin for worn out, then the Latin for wakeful, then say \emph{tēgula, tēgulae} again
\end{itemize}

\begin{tcolorbox}[breakable,colback=teal!4,colframe=teal!35!black,title={Before you move on}]
What does tēgula, tēgulae mean? (\textquotedblleft{}A roof tile\textquotedblright{}.) Say it once more.
\end{tcolorbox}
\section[fornāx, fornācis]{fornāx, fornācis — a furnace, a kiln}

\label{lesson:LA-C217-fornax}



\begin{quote}
Before the new one: say the Latin for wakeful, then the Latin for a roof tile.
\end{quote}

\subsection*{You'll want to know: fornāx, fornācis}
\textbf{fornāx, fornācis} — \textquotedblleft{}a furnace, a kiln\textquotedblright{}.

\begin{cousinweb}
\textbf{fornāx} belongs with \emph{fornus}, \emph{furnus}, \textquotedblleft{}an oven\textquotedblright{} (Lewis \& Short), usually traced to Proto-Indo-European \emph{*gwher-}, \textquotedblleft{}warm\textquotedblright{}; English \emph{warm} is a cousin.
\end{cousinweb}

\begin{grammarlens}[title={What you've built}]
One more word for this chapter.
\end{grammarlens}

\subsection*{Your turn}
\begin{itemize}
\raggedright
  \item \emph{Say it:} \emph{fornāx, fornācis}
  \item \emph{Say it:} \emph{fornāx, fornācis}, once more
  \item \emph{Say it:} say \emph{fornāx, fornācis}
  \item \emph{Recall:} say the Latin for wakeful, then the Latin for a roof tile, then say \emph{fornāx, fornācis} again
\end{itemize}

\begin{tcolorbox}[breakable,colback=teal!4,colframe=teal!35!black,title={Before you move on}]
What does fornāx, fornācis mean? (\textquotedblleft{}A furnace, a kiln\textquotedblright{}.) Say it once more.
\end{tcolorbox}
\section[lectīca, lectīcae]{lectīca, lectīcae — a litter (carried chair)}

\label{lesson:LA-C217-lectica}



\begin{quote}
Before the new one: say the Latin for a roof tile, then the Latin for a furnace.
\end{quote}

\subsection*{You'll want to know: lectīca, lectīcae}
\textbf{lectīca, lectīcae} — \textquotedblleft{}a litter (carried chair)\textquotedblright{}.

\begin{cousinweb}
\textbf{lectīca} is built on \emph{lectus}, \textquotedblleft{}a bed, a couch\textquotedblright{} (EtymDB 2.1): a couch carried on poles.
\end{cousinweb}

\begin{grammarlens}[title={What you've built}]
One more word for this chapter.
\end{grammarlens}

\subsection*{Your turn}
\begin{itemize}
\raggedright
  \item \emph{Say it:} \emph{lectīca, lectīcae}
  \item \emph{Say it:} \emph{lectīca, lectīcae}, once more
  \item \emph{Say it:} say \emph{lectīca, lectīcae}
  \item \emph{Recall:} say the Latin for a roof tile, then the Latin for a furnace, then say \emph{lectīca, lectīcae} again
\end{itemize}

\begin{tcolorbox}[breakable,colback=teal!4,colframe=teal!35!black,title={Before you move on}]
What does lectīca, lectīcae mean? (\textquotedblleft{}A litter (carried chair)\textquotedblright{}.) Say it once more.
\end{tcolorbox}
\section[vehiculum, vehiculī]{vehiculum, vehiculī — a vehicle}

\label{lesson:LA-C217-vehiculum}



\begin{quote}
Before the new one: say the Latin for a furnace, then the Latin for a litter.
\end{quote}

\subsection*{You'll want to know: vehiculum, vehiculī}
\textbf{vehiculum, vehiculī} — \textquotedblleft{}a vehicle\textquotedblright{}.

\begin{cousinweb}
\textbf{vehiculum} is built on \emph{vehō}, \textquotedblleft{}to carry\textquotedblright{} (Lewis \& Short), with \emph{-culum}, an ending for a means: a means of carrying.
\end{cousinweb}

\begin{grammarlens}[title={What you've built}]
One more word for this chapter.
\end{grammarlens}

\subsection*{Your turn}
\begin{itemize}
\raggedright
  \item \emph{Say it:} \emph{vehiculum, vehiculī}
  \item \emph{Say it:} \emph{vehiculum, vehiculī}, once more
  \item \emph{Say it:} say \emph{vehiculum, vehiculī}
  \item \emph{Recall:} say the Latin for a furnace, then the Latin for a litter, then say \emph{vehiculum, vehiculī} again
\end{itemize}

\begin{tcolorbox}[breakable,colback=teal!4,colframe=teal!35!black,title={Before you move on}]
What does vehiculum, vehiculī mean? (\textquotedblleft{}A vehicle\textquotedblright{}.) Say it once more.
\end{tcolorbox}
\section[currus, currūs]{currus, currūs — a chariot}

\label{lesson:LA-C217-currus}



\begin{quote}
Before the new one: say the Latin for a litter, then the Latin for a vehicle.
\end{quote}

\subsection*{You'll want to know: currus, currūs}
\textbf{currus, currūs} — \textquotedblleft{}a chariot\textquotedblright{}.

\begin{cousinweb}
\textbf{currus} is built on \emph{currō}, \textquotedblleft{}to run\textquotedblright{} (Lewis \& Short), usually traced to Proto-Indo-European \emph{*kers-}, \textquotedblleft{}to run\textquotedblright{} (EtymDB 2.1). English \emph{car} comes from a Celtic cousin of the same root, through Latin \emph{carrus}, \textquotedblleft{}a wagon\textquotedblright{}.
\end{cousinweb}

\begin{grammarlens}[title={What you've built}]
One more word for this chapter.
\end{grammarlens}

\subsection*{Your turn}
\begin{itemize}
\raggedright
  \item \emph{Say it:} \emph{currus, currūs}
  \item \emph{Say it:} \emph{currus, currūs}, once more
  \item \emph{Say it:} say \emph{currus, currūs}
  \item \emph{Recall:} say the Latin for a litter, then the Latin for a vehicle, then say \emph{currus, currūs} again
\end{itemize}

\begin{tcolorbox}[breakable,colback=teal!4,colframe=teal!35!black,title={Before you move on}]
What does currus, currūs mean? (\textquotedblleft{}A chariot\textquotedblright{}.) Say it once more.
\end{tcolorbox}
